% ===== Addiction: "Dealing with the monkey on your back" (~17 min) =====
\partaccent{mint}
\topicopener{monkey}{Կախվածություն}{Կապիկը մեջքիդ}{Ինչպես է ուղեղը թմրանյութից կախվածություն ձևավորում}{hungry}

\begin{frame}{ԱՄՆ, 1987. գովազդ ձվով}
  \centering
  % The 1987 US public-service ad, redrawn: egg = brain, frying pan = drugs, fried egg = brain on drugs.
% Panels reveal on slides 1-3.
\begin{tikzpicture}[pcap/.style={font=\blackfont\small, text=plum, align=center, text width=3.8cm}]
  % 1. the egg
  \filldraw[fill=cardwhite, draw=plum, line width=1.2pt]
    (0,-0.75) .. controls (0.85,-0.75) and (0.8,0.55) .. (0,0.85) .. controls (-0.8,0.55) and (-0.85,-0.75) .. (0,-0.75);
  \node[pcap] at (0,-1.45) {«Սա քո ուղեղն է»};
  % 2. the pan
  \begin{scope}[shift={(4.3,0)}, visible on=<2->]
    \fill[plum!85!black] (0.85,-0.1) rectangle (2.0,0.1);
    \filldraw[fill=plum!75!cardwhite, draw=plum, line width=1.2pt] (0,0) circle (0.95);
    \fill[plum!55!cardwhite] (0,0) circle (0.78);
    \node[pcap] at (0.4,-1.45) {«Սա թմրանյութն է»};
  \end{scope}
  % 3. the fried egg
  \begin{scope}[shift={(9.0,0)}, visible on=<3->]
    \fill[plum!85!black] (0.85,-0.1) rectangle (2.0,0.1);
    \filldraw[fill=plum!75!cardwhite, draw=plum, line width=1.2pt] (0,0) circle (0.95);
    \filldraw[fill=cardwhite, draw=plum!50, line width=0.6pt]
      plot[smooth cycle, tension=0.8] coordinates {(-0.62,0.05) (-0.35,0.55) (0.2,0.62) (0.6,0.3) (0.58,-0.3) (0.15,-0.62) (-0.4,-0.5)};
    \filldraw[fill=amOrange, draw=amOrange!80!black, line width=0.6pt] (-0.02,0.02) circle (0.24);
    \node[pcap] at (0.4,-1.45) {«Սա քո ուղեղն է\\ թմրանյութերից»};
  \end{scope}
\end{tikzpicture}
\par\bigskip
  \raggedright
  \onslide<4->{%
  \begin{poll}
    Գովազդը մրցանակներ ստացավ։ Իսկ գիտականորեն ճի՞շտ է։ Ռեակցիա. \emoji{thumbsup} կամ \emoji{thumbsdown}
  \end{poll}}
\end{frame}

\begin{frame}{Լավ գովազդ, վատ գիտություն}
  \begin{columns}[c]
    \begin{column}{0.56\textwidth}
      \begin{enumerate}
        \item Թմրանյութն ուղեղն այնքան չի տաքացնում, որ սպիտակուցները «եփվեն»
        \item<2-> Հազվադեպ է ազդում \textbf{ամբողջ} ուղեղի վրա միանգամից, ինչպես տապակը՝ ձվի
        \item<3-> Եվ ուղեղը «կճեպից» (գանգից) չենք հանում։ Այլապես թմրանյութերն այդքան տարածված չէին լինի
      \end{enumerate}
    \end{column}
    \begin{column}{0.38\textwidth}
      \onslide<4->{%
      \begin{bookquote}
        Սա չի նշանակում, որ թմրանյութերն օգտակար են ուղեղի համար։ Պարզապես ճշմարտությունը շատ ավելի բարդ է, քան ձվով փոխաբերությունները։
      \end{bookquote}}
    \end{column}
  \end{columns}
\end{frame}

\begin{frame}{Վտանգավոր են, որովհետև աշխատում են}
  \begin{columns}[c]
    \begin{column}{0.4\textwidth}
      \begin{card}[colback=butter!45!cream]
        \centering\bigemoji[1cm]{moneybag}\par\smallskip
        Ապօրինի թմրանյութերի շուկան, գնահատականներով, գրեթե \textbf{կես տրիլիոն դոլար} է
      \end{card}
    \end{column}
    \begin{column}{0.54\textwidth}
      \begin{itemize}
        \item Ասում ենք՝ թմրանյութերը վնասում են առողջությանն ու կործանում կյանքեր։ Ու դա ճիշտ է
        \item<2-> Բայց դա լինում է, որովհետև դրանք \textbf{շատ լավ են աշխատում}. փոխում են ուղեղի ամենահիմնական պրոցեսները
        \item<3-> Արդյունքը՝ կախվածություն, վարքի փոփոխություններ ու ավելին
      \end{itemize}
    \end{column}
  \end{columns}
\end{frame}

\begin{frame}{Պարգևի ուղին}
  \begin{columns}[c]
    \begin{column}{0.46\textwidth}
      \centering
      % The reward pathway on the medial brain drawing (front = left; images/brain_medial.png, P. J. Lynch, CC BY 2.5).
% Positions are approximate fractions of the image (0,0 = bottom left). Slide 2 adds the link to the frontal cortex.
\begin{tikzpicture}
  \node[anchor=south west, inner sep=0] (b) at (0,0) {\includegraphics[width=5cm]{brain_medial.png}};
  \begin{scope}[x={(b.south east)}, y={(b.north west)}]
    \node[braindot=amRed] (vta) at (0.49,0.37) {};
    \node[braindot=amRed] (nac) at (0.35,0.47) {};
    \draw[arr=amRed, line width=1.6pt, shorten >=1mm, shorten <=1mm] (vta) to[bend left=20] (nac);
    \node[brainlbl=brainpink!40!cardwhite, anchor=north west] (vl) at (0.47,-0.03) {VTA, դոֆամինի «աղբյուր»};
    \draw[plum, line width=0.6pt] (vta) -- (vl.north west);
    \node[brainlbl=brainpink!40!cardwhite, anchor=north east] (nl) at (0.38,-0.03) {nucleus accumbens};
    \draw[plum, line width=0.6pt] (nac) -- (nl.north east);
    \node[braindot=amBlue, visible on=<2->] (pfc) at (0.12,0.58) {};
    \draw[arr=amBlue, line width=1.4pt, dashed, shorten >=1mm, shorten <=1mm, visible on=<2->] (nac) to[bend left=25] (pfc);
    \node[brainlbl=lavender!50!cardwhite, anchor=south west, visible on=<2->] (pl) at (0.0,1.02) {նախաճակատային կեղև («արգելակ»)};
    \draw[plum, line width=0.6pt, visible on=<2->] (pfc) -- (pl.south west -| pfc);
  \end{scope}
\end{tikzpicture}

    \end{column}
    \begin{column}{0.5\textwidth}
      \begin{itemize}
        \item Դոֆամինային \textbf{մեզոլիմբիկ ուղին} հաճելի զգացումով «պարգևատրում» է օգտակար արարքների համար
        \item Համեղ մանդարինից մինչև ննջասենյակային որոշ զբաղմունքներ
        \item Շաքարն արագ էներգիա է, դրա համար էլ քաղցրը հաճելի է
      \end{itemize}
      \onslide<2->{%
      \begin{card}[colback=mint!30!cardwhite]
        Ամիսներ ծովում մոլորվածի համար ջուրն ու հացը \textbf{դրախտային} կերակուր են
      \end{card}}
    \end{column}
  \end{columns}
\end{frame}

\begin{frame}{Կարճ ճանապարհ դեպի պարգև}
  \centering
  \begin{tikzpicture}[st/.style={flow=2.6cm, font=\small}]
    \node[st, fill=butter!50!cardwhite] (f) at (0,0.9) {\emoji{bread}\ սնունդ, ջուր};
    \node[st, fill=mint!40!cardwhite] (b) at (3.4,0.9) {օգուտ\\ մարմնին};
    \node[st, fill=lavender!40!cardwhite] (r) at (6.8,0.9) {ուղեղը\\ «ճանաչում» է};
    \node[st, fill=brainpink!45!cardwhite] (p) at (10.2,0) {\emoji{party}\\ \textbf{պարգև}};
    \draw[arr] (f) -- (b); \draw[arr] (b) -- (r); \draw[arr] (r) -- (p);
    \node[st, fill=peach!50!cardwhite, visible on=<2->] (d) at (0,-0.9) {\emoji{pill}\ թմրանյութ};
    \draw[arr=pinkdark, line width=2pt, visible on=<2->] (d) -- (p);
    \node[font=\blackfont\small, text=pinkdark, visible on=<2->] at (5.6,-1.05) {ուղիղ, առանց «մանրամասների»};
  \end{tikzpicture}\par\bigskip
  \raggedright
  \onslide<3->{%
  \begin{bookquote}
    Ինչպես բանկի աշխատակից, որ փողով լի պայուսակներ է տալիս՝ առանց հաշվեհամարի ու անձնագրի նման ձանձրալի մանրուքների։
  \end{bookquote}}
\end{frame}

\begin{frame}{Բանալիներ ու կողպեքներ}
  \begin{columns}[T]
    \begin{column}{0.31\textwidth}
      \begin{card}[colback=mint!35!cardwhite]
        \emoji[1.3em]{key}\emoji[1.3em]{lock}\ {\blackfont Նեյրոմեդիատոր}\par\smallskip
        Ճիշտ բանալին բացում է իր կողպեքը (հիշո՞ւմ ես սինապսը)
      \end{card}
    \end{column}
    \begin{column}{0.31\textwidth}
      \onslide<2->{%
      \begin{card}[colback=lavender!40!cardwhite]
        \emoji[1.3em]{key}\ {\blackfont Ագոնիստ}\par\smallskip
        Այլ նյութ, որ \textbf{նույնպես} բացում է կողպեքը։ Օր.՝ սրտի դեղեր, որ նմանակում են ադրենալինը
      \end{card}}
    \end{column}
    \begin{column}{0.31\textwidth}
      \onslide<3->{%
      \begin{card}[colback=peach!40!cardwhite]
        \emoji[1.3em]{stop}\ {\blackfont Անտագոնիստ}\par\smallskip
        Զբաղեցնում է ընկալիչը, բայց չի բացում, ինչպես վերելակի դռների մեջ խրված ճամպրուկը
      \end{card}}
    \end{column}
  \end{columns}
  \vspace{3mm}
  \onslide<4->{%
  \begin{card}[colback=butter!45!cream]
    Հակափսիխոտիկ դեղերը \textbf{փակում} են դոֆամինի որոշ ընկալիչներ, իսկ թմրանյութերն արհեստականորեն \textbf{միացնում} են պարգևի ուղին
  \end{card}}
\end{frame}

\begin{frame}{Բոլոր ճանապարհները տանում են դոֆամինի}
  \begin{columns}[c]
    \begin{column}{0.5\textwidth}
      {\small Պարգևի ուղին ունի շատ տարբեր ընկալիչներ, ու շատ նյութեր կարող են այն գրգռել.}\par\smallskip
      \chip{}{կոկաին}\ \chip{}{հերոին}\ \chip{cigarette}{նիկոտին}\par\smallskip
      \chip{}{ամֆետամիններ}\ \chip{wine}{ալկոհոլ}\par\medskip
      \onslide<2->{Բոլորն էլ ավելացնում են \textbf{դոֆամինի} հաղորդումը պարգևի ուղում։ Դրա համար են «հաճելի»}
    \end{column}
    \begin{column}{0.46\textwidth}
      \onslide<3->{%
      \begin{scienceupdate}
        Գրքում կոկաինը «նմանակում է դոֆամինը»։ Իրականում այն \textbf{փակում է դոֆամինի «պոմպը»} (transporter), ու դոֆամինը կուտակվում է սինապսում։ Ինչպես SSRI-ն՝ սերոտոնինի դեպքում
      \end{scienceupdate}}
    \end{column}
  \end{columns}
  \only<3->{\source{NIDA, Cocaine Research Report}}
\end{frame}

\begin{frame}{Դոֆամինը «հաճույքի հորմոն» չէ}
  \begin{columns}[c]
    \begin{column}{0.62\textwidth}
      \begin{scienceupdate}
        Դոֆամինն ավելի շատ \textbf{«ուզել»}-ի (wanting), մոտիվացիայի համար է, քան բուն հաճույքի։ Հաճույքը («հավանել», liking) կախված է այլ համակարգերից, օրինակ՝ ուղեղի սեփական ափիոնային նյութերից։\par\smallskip
        Սա ցույց են տվել Ռոբինսոնն ու Բերիջը. գիրքը նրանց աշխատանքին հղում է նաև կախվածության բաժնում։
      \end{scienceupdate}
    \end{column}
    \begin{column}{0.34\textwidth}
      \onslide<2->{%
      \begin{card}[colback=lavender!35!cardwhite]
        \centering\bigemoji[1cm]{monkey}\par\smallskip
        \raggedright Կախվածության ժամանակ կարող ես շատ \textbf{ուզել} մի բան, որն արդեն \textbf{չես վայելում}
      \end{card}}
    \end{column}
  \end{columns}
  \source{Berridge \& Robinson, Brain Res. Rev., 1998}
\end{frame}

\begin{frame}{Արգելակները}
  \begin{columns}[c]
    \begin{column}{0.54\textwidth}
      \begin{itemize}
        \item Ուղեղն արագ է հասկանում, թե ինչն է հաճելի, ու ինչպես ստանալ ավելին
        \item Բայց կան նաև «արգելակներ»՝ ինքնատիրապետման կենտրոններ, ամենայն հավանականությամբ՝ \textbf{նախաճակատային կեղևում}
        \item<3-> Դրանք թույլ չեն տալիս ամբողջովին տրվել հաճույքներին
      \end{itemize}
    \end{column}
    \begin{column}{0.4\textwidth}
      \centering
      \onslide<2->{%
      \begin{tikzpicture}
        \node at (0,0) {\mascot[2.2cm]{hungry}};
        \node[say, text width=3.2cm, font=\handfont\large, callout absolute pointer={(0.3,0.9)}] at (0.6,2.3) {Բան լավ։\\ Ավելի բան պետք։};
      \end{tikzpicture}}
    \end{column}
  \end{columns}
\end{frame}

\begin{frame}{Տոլերանտություն. քաղաք, ուր համերգ է եկել}
  \begin{columns}[c]
    \begin{column}{0.5\textwidth}
      \begin{card}[colback=butter!45!cream]
        \emoji[1.3em]{guitar}\ Հազարավոր ոգևորված երկրպագուներ են գալիս քաղաք, ու քաոս է սկսվում
      \end{card}
      \onslide<2->{%
      \begin{card}[colback=lavender!40!cardwhite]
        \emoji[1.3em]{police}\ Քաղաքը պատասխանում է. ավելի շատ ոստիկան, փակ փողոցներ, ավելի հաճախ ավտոբուսներ
      \end{card}}
    \end{column}
    \begin{column}{0.44\textwidth}
      \onslide<3->{%
      \begin{itemize}
        \item Երկրպագուները թմրանյութն են, քաղաքը՝ ուղեղը
        \item Ուղեղը ճնշում է ընկալիչները կամ «անջատում» դրանք, որ ակտիվությունը նորմալ դառնա
        \item Սա \textbf{տոլերանտությունն} է (tolerance). նույն դոզան այլևս նույն ազդեցությունը չունի
      \end{itemize}}
    \end{column}
  \end{columns}
\end{frame}

\begin{frame}{Ավելի ու ավելի մեծ դոզա}
  \begin{columns}[c]
    \begin{column}{0.48\textwidth}
      \centering
      % Schematic tolerance staircase: the dose needed for the same effect keeps growing, until it would kill a first-timer.
% Not data: columns of pill icons, revealed on slides 1-4; the red line comes on slide 5.
\begin{tikzpicture}
  \foreach \c/\n/\s/\t in {0/1/1/սկիզբ, 1/2/2/հետո, 2/4/3/ավելի ուշ, 3/6/4/շատ ավելի ուշ}{
    \foreach \k in {1,...,\n}{\node[visible on=<\s->] at (\c*1.55,\k*0.52-0.2) {\bigemoji[0.46cm]{pill}};}
    \node[font=\scriptsize, text=plum, align=center, text width=1.45cm, visible on=<\s->] at (\c*1.55,-0.35) {\t};
  }
  \draw[plum, line width=0.8pt] (-0.7,-0.08) -- (5.4,-0.08);
  \draw[amRed, dashed, line width=1.4pt, visible on=<5->] (-0.7,2.75) -- (5.4,2.75);
  \node[anchor=south west, font=\blackfont\scriptsize, text=amRed, visible on=<5->, inner sep=1pt] at (-0.7,2.8) {նորեկին կսպաներ};
\end{tikzpicture}
\par
      {\scriptsize սխեմատիկ}
    \end{column}
    \begin{column}{0.48\textwidth}
      \begin{itemize}
        \item Ուղեղը հարմարվեց. նույն զգացողության համար ավելի մեծ դոզա է պետք
        \item<2-> Ուղեղը նորից հարմարվում է, ու այդպես շարունակ
        \item<5-> Վերջում մարդն ընդունում է դոզաներ, որոնք \textbf{կսպանեին} առաջին անգամ փորձողին, ու ստանում է նույն զգացողությունը, ինչ սկզբում
      \end{itemize}
    \end{column}
  \end{columns}
\end{frame}

\begin{frame}{Ինչու է այդքան դժվար թողնելը}
  \begin{columns}[c]
    \begin{column}{0.46\textwidth}
      \begin{itemize}
        \item Միայն կամքի հարց չէ. ուղեղն ու մարմինը \textbf{ֆիզիկապես} փոխվել են
        \item Ափիոնային նյութերը (օր.՝ հերոին) ճնշում են ցավը և էյֆորիա են տալիս
        \item<2-> Ցավը կարևոր ազդանշան է, ու ուղեղը \textbf{ուժեղացնում է} ցավի համակարգը
        \item<3-> Կտրուկ թողնելիս (cold turkey) մնում է միայն ուժեղացված ցավը։ Ծանր է ու վտանգավոր
      \end{itemize}
    \end{column}
    \begin{column}{0.5\textwidth}
      \centering
      % Opiates vs the pain system as a seesaw (schematic, after the book):
% slide 1 the opiate wins (euphoria), slide 2 the brain boosts the pain system (balance), slide 3 the drug is gone (withdrawal).
\begin{tikzpicture}[
    wt/.style={draw=plum, line width=0.8pt, rounded corners=1.5mm, font=\small, text=plum, align=center, inner sep=3pt},
  ]
  \filldraw[fill=plum, draw=plum] (-0.35,-1.1) -- (0.35,-1.1) -- (0,-0.2) -- cycle;
  \only<1>{\begin{scope}[rotate around={8:(0,-0.2)}]
    \draw[plum, line width=3pt, line cap=round] (-2.9,-0.2) -- (2.9,-0.2);
    \node[wt, fill=mint!50!cardwhite, anchor=south, minimum width=1.9cm, minimum height=1.2cm] at (-2.0,-0.12) {\emoji{pill}\\ ափիոնային\\ նյութ};
    \node[wt, fill=brainpink!40!cardwhite, anchor=south, minimum width=1.3cm, minimum height=0.6cm] at (2.0,-0.12) {ցավ};
  \end{scope}}
  \only<2>{\begin{scope}
    \draw[plum, line width=3pt, line cap=round] (-2.9,-0.2) -- (2.9,-0.2);
    \node[wt, fill=mint!50!cardwhite, anchor=south, minimum width=1.9cm, minimum height=1.2cm] at (-2.0,-0.12) {\emoji{pill}\emoji{pill}\\ ավելի շատ\\ ափիոն};
    \node[wt, fill=brainpink!60!cardwhite, anchor=south, minimum width=1.9cm, minimum height=1.2cm] at (2.0,-0.12) {\emoji{lightning}\\ ուժեղացված\\ ցավի համակարգ};
  \end{scope}}
  \only<3->{\begin{scope}[rotate around={-16:(0,-0.2)}]
    \draw[plum, line width=3pt, line cap=round] (-2.9,-0.2) -- (2.9,-0.2);
    \node[wt, fill=amRed!30!cardwhite, anchor=south, minimum width=1.9cm, minimum height=1.2cm] at (2.0,-0.12) {\emoji{lightning}\emoji{lightning}\\ ցավ\\ ամեն ինչից};
  \end{scope}
  \node[font=\blackfont\small, text=pinkdark] at (-2.2,1.6) {դեղը չկա};}
\end{tikzpicture}

    \end{column}
  \end{columns}
\end{frame}

\begin{frame}{Ամեն ինչ, ինչ կապված է թմրանյութի հետ}
  \begin{columns}[c]
    \begin{column}{0.5\textwidth}
      \begin{itemize}
        \item \textbf{Սենսիտիզացիա}՝ տոլերանտության հակառակը. թմրանյութն ավելի ու ավելի ուժեղ է «կանչում»
        \item<2-> Պարգևի ուղին ու \textbf{նշաձև մարմինը} զգացմունքային «լիցք» են տալիս այն ամենին, ինչ կապված է թմրանյութի հետ (drug cues)
        \item<3-> Ծխամորճը, ներարկիչը, կրակայրիչը, հոտը. դրանք արդեն ինքնին գրգռում են
      \end{itemize}
    \end{column}
    \begin{column}{0.44\textwidth}
      \onslide<4->{%
      \begin{card}[colback=peach!35!cardwhite]
        \emoji[1.3em]{syringe}\ {\blackfont Մեթադոնի պատմությունը}\par\smallskip
        Մեթադոնն օգնում է աստիճանաբար թողնել հերոինը. այն խմում են։ Բայց որոշ մարդիկ այն թքում են ներարկիչի մեջ ու \textbf{ներարկում}, որովհետև հենց ներարկելու արարքն է «հաճելի» դարձել
      \end{card}}
    \end{column}
  \end{columns}
\end{frame}

\begin{frame}{Երբ «արգելակները» թուլանում են}
  \begin{columns}[c]
    \begin{column}{0.54\textwidth}
      \begin{itemize}
        \item Պարգևի ուղու ու ճակատային կեղևի կապը փոխվում է. թմրանյութը դառնում է ավելի կարևոր, քան աշխատանքը, օրենքը, ցնցուղ ընդունելը
        \item<2-> Իսկ բացասական հետևանքները (ձերբակալություն, հիվանդություն, ընկերների կորուստ) քիչ են անհանգստացնում
        \item<3-> Նախաճակատային կեղևի «Մի՛ արա», «Կփոշմանես» ձայները թուլանում են
      \end{itemize}
    \end{column}
    \begin{column}{0.4\textwidth}
      \onslide<4->{%
      \begin{bookquote}
        Ազատ կամքը կարող է մարդու ուղեղի ամենամեծ նվաճումներից մեկը լինել, բայց եթե այն խանգարում է հաճույքին, ուրեմն պետք է գնա։
      \end{bookquote}}
    \end{column}
  \end{columns}
\end{frame}

\begin{frame}{Ինչու ոմանք, բայց ոչ բոլորը}
  \begin{columns}[c]
    \begin{column}{0.52\textwidth}
      \begin{itemize}
        \item Մռայլ թաղամասեր, որտեղ թմրանյութը միակ «փախուստն» է
        \item Չախտորոշված խանգարում ու «ինքնաբուժում»
        \item<2-> Գեներ. գուցե ինքնատիրապետման կենտրոններն են թույլ
      \end{itemize}
      \onslide<3->{%
      \begin{card}[colback=brainpink!25!cardwhite]
        \emoji{warning}\ Թողնելուց հետո փոփոխությունները չեն անհետանում, պարզապես «չեն օգտագործվում»։ Դրա համար է \textbf{ռեցիդիվը} այդքան հեշտ
      \end{card}}
    \end{column}
    \begin{column}{0.42\textwidth}
      \onslide<2->{%
      \begin{tikzpicture}
        \node[say, text width=3.6cm, font=\handfont\large, callout relative pointer={(-0.3,-0.5)}] at (0,1.3) {«Ի՞նչ վատ բան կարող է լինել որ...»};
        \node[say, text width=3.6cm, font=\handfont\large, fill=lavender!35!cardwhite, callout relative pointer={(0.3,-0.5)}, visible on=<3->] at (0.5,-0.5) {«Ահա ամբողջական ցուցակը, մանրամասն...»};
      \end{tikzpicture}\par\smallskip
      \onslide<3->{\small Երկրորդ ձայնը ոմանց մոտ շատ թույլ է}}
    \end{column}
  \end{columns}
\end{frame}

\begin{frame}{Ոչ միայն թմրանյութեր}
  \begin{columns}[c]
    \begin{column}{0.48\textwidth}
      \begin{chatbox}
        Ուրիշ ինչի՞ց կարող է մարդը կախվածություն ունենալ։ Գրիր չաթում։
      \end{chatbox}
      \onslide<2->{%
      \centering\bigemoji[0.9cm]{videogame}\ \bigemoji[0.9cm]{moneybag}\ \bigemoji[0.9cm]{slotmachine}\par
      {\small գնումներ, համակարգչային խաղեր, մոլախաղեր}}
    \end{column}
    \begin{column}{0.48\textwidth}
      \onslide<3->{%
      \begin{card}[colback=peach!35!cardwhite]
        \emoji[1.3em]{slotmachine}\ {\blackfont Մոլախաղը շատ դժվար է թողնել}\par\smallskip
        Կախվածությունը «ջնջելու» համար ուղեղը պետք է երկար ժամանակ պարգև չստանա։ Բայց մոլախաղում երկար պարտվելը \textbf{նորմալ} է, ու ուղեղը շարունակում է սպասել շահումին
      \end{card}}
    \end{column}
  \end{columns}
\end{frame}

\begin{frame}{Բուժելի վիճակ, ոչ թե բնավորության թերություն}
  \begin{columns}[c]
    \begin{column}{0.56\textwidth}
      \begin{itemize}
        \item Ինչպես էլ որ սկսվել է, մասնագետները կախվածությունը համարում են վիճակ, որը պետք է \textbf{բուժել}, ոչ թե դատապարտել
        \item<2-> Թմրանյութն ուղեղը շրջում է ինքն իր դեմ՝ երկար պատերազմում, որի մարտադաշտը մեր կյանքն է
      \end{itemize}
    \end{column}
    \begin{column}{0.38\textwidth}
      \onslide<3->{%
      \begin{bookquote}
        Սա սարսափելի բան է՝ ինքդ քեզ անելու համար։ Բայց թմրանյութերն անում են այնպես, որ քեզ համար միևնույն լինի։
      \end{bookquote}}
    \end{column}
  \end{columns}
\end{frame}

\begin{frame}{Ձվերով դժվար է բացատրել}
  \centering
  \begin{meme}[7.2cm]{tuxedo_pooh.png}
    \memecap[text width=3.3cm]{0.72}{0.75}{«Սա քո ուղեղն է թմրանյութերից» \emoji{fire}}
    \memecap[text width=3.6cm, font=\blackfont\scriptsize]{0.72}{0.25}{Դոֆամին, ընկալիչների հարմարում, տոլերանտություն, սենսիտիզացիա, նշաձև մարմին, նախաճակատային կեղև}
  \end{meme}
\end{frame}
